% ============================================================
%  第5章：凸集与凸函数
% ============================================================
\section{第5章 \quad 凸集与凸函数}

\subsection{5.1 为什么非线性规划需要凸性}

线性规划天然是"好的"——可行域是凸多面体，局部最优就是全局最优，单纯形法保证找到全局最优解。但到了非线性规划，目标函数与可行域都可能非凸，局部最优不再自动等于全局最优。

凸性是判断"局部 $\Rightarrow$ 全局"是否成立的关键条件。本章构成非线性规划的理论基础，直接服务于第6章（凸规划下KKT条件变为充要条件）和第8章（共轭梯度法/拟牛顿法的收敛性分析依赖凸性假设）。

\textbf{核心逻辑链：} 凸集 $\to$ 凸函数 $\to$ 凸规划 $\to$ 局部极小=全局极小 $\to$ KKT条件充要。

\textbf{考试定位：} 凸性判别嵌入于非线性规划的各类问题中。判定Hessian正定是第8章（Newton法/拟Newton法）的必备步骤。判断一个问题是否为凸规划，直接影响KKT点是局部最优还是全局最优的结论。

% ============================================================
\subsection{5.1A 预备知识：多元微积分基础}

\subsubsection{偏导数}

对于 $n$ 元函数 $f(x_1,\dots,x_n)$，偏导数 $\frac{\partial f}{\partial x_i}$ 表示将其他变量视为常数时 $f$ 对 $x_i$ 的导数。

\subsubsection{梯度 $\nabla f(x)$ —— 一阶导数信息}

梯度是 $f$ 的所有一阶偏导数排成的列向量：
\[
\nabla f(x) = \begin{pmatrix}
\frac{\partial f}{\partial x_1} \\
\frac{\partial f}{\partial x_2} \\
\vdots \\
\frac{\partial f}{\partial x_n}
\end{pmatrix}
\]

\textbf{二次函数：} $f(x)=\frac12 x^T Q x + c^T x$（$Q$ 对称），则 $\nabla f(x)=Qx+c$。

\subsubsection{Hessian 矩阵 $\nabla^2 f(x)$ —— 二阶导数信息}

Hessian 是 $f$ 的所有二阶偏导数排成的 $n\times n$ 对称矩阵：
\[
[\nabla^2 f(x)]_{ij} = \frac{\partial^2 f}{\partial x_i \partial x_j}(x)
\]

\[
\nabla^2 f(x) = \begin{pmatrix}
\frac{\partial^2 f}{\partial x_1^2} & \frac{\partial^2 f}{\partial x_1\partial x_2} & \cdots & \frac{\partial^2 f}{\partial x_1\partial x_n} \\[4pt]
\frac{\partial^2 f}{\partial x_2\partial x_1} & \frac{\partial^2 f}{\partial x_2^2} & \cdots & \frac{\partial^2 f}{\partial x_2\partial x_n} \\
\vdots & \vdots & \ddots & \vdots \\
\frac{\partial^2 f}{\partial x_n\partial x_1} & \frac{\partial^2 f}{\partial x_n\partial x_2} & \cdots & \frac{\partial^2 f}{\partial x_n^2}
\end{pmatrix}
\]

\textbf{二次函数的Hessian：} $f(x)=\frac12 x^T Q x + c^T x$（$Q$ 对称），$\nabla^2 f(x)=Q$（常数矩阵）。这就是第8章牛顿法对正定二次函数一步到位的原因——Hessian不随 $x$ 变化，二次模型就是函数本身。

\subsubsection{常用导数/梯度公式}

\begin{itemize}
  \item $\nabla(b^T x) = b$（$b$ 为常向量）
  \item $\nabla(x^T A x) = 2Ax$（$A$ 对称）；若 $A$ 不对称则为 $(A+A^T)x$
  \item $\nabla(\frac12 x^T Q x + c^T x) = Qx + c$（$Q$ 对称）
  \item $\nabla\|x\|_2^2 = \nabla(x^T x) = 2x$
\end{itemize}

\subsubsection{Taylor 展开}

\textbf{一阶Taylor展开：} $f(y) \approx f(x) + \nabla f(x)^T(y-x)$。右边是关于 $y$ 的线性函数——在 $x$ 处用切平面代替原函数。

\textbf{二阶Taylor展开（令 $d=y-x$）：}
\[
f(x+d) \approx f(x) + \nabla f(x)^T d + \frac{1}{2} d^T \nabla^2 f(x) d
\]
右边是关于 $d$ 的二次函数——在 $x$ 处用二次曲面（抛物面）代替原函数。牛顿法就是直接跳到这个二次逼近的极小点。

\subsubsection{方向导数}

$f$ 在 $x$ 处沿方向 $d$ 的方向导数（directional derivative）是一阶变化率：
\[
D_d f(x) = \lim_{t\to 0^+} \frac{f(x+td)-f(x)}{t} = \nabla f(x)^T d
\]

\textbf{第6章所有最优性条件几何解释的根基：}
\begin{itemize}
  \item $\nabla f(x)^T d < 0$ $\Longleftrightarrow$ 沿 $d$ 方向走一小步，$f$ 的值下降
  \item $\nabla f(x)^T d = 0$ $\Longleftrightarrow$ 沿 $d$ 方向走一小步，$f$ 的值不变（一阶意义上）
  \item $\nabla f(x)^T d > 0$ $\Longleftrightarrow$ 沿 $d$ 方向走一小步，$f$ 的值上升
\end{itemize}

\subsubsection{矩阵正定/半正定/负定的形式定义}

设 $A$ 为 $n\times n$ 实对称矩阵（$A^T=A$）：
\begin{itemize}
  \item $A$ \textbf{正定}（$A \succ 0$）：$\forall x\neq 0$，$x^T A x > 0$
  \item $A$ \textbf{半正定}（$A \succeq 0$）：$\forall x$，$x^T A x \ge 0$
  \item $A$ \textbf{负定}（$A \prec 0$）：$\forall x\neq 0$，$x^T A x < 0$（等价于 $-A \succ 0$）
  \item $A$ \textbf{不定}：存在 $x,y$ 使得 $x^T A x > 0$ 而 $y^T A y < 0$
\end{itemize}

形式定义与顺序主子式法的关系：$x^T A x > 0$（$\forall x\neq0$）是正定的定义；顺序主子式法是判定方法（Sylvester判别法）——两者等价。但半正定的顺序主子式判别法有陷阱，详见下文。

\subsubsection{$\operatorname{diag}$ 记号}

$\operatorname{diag}(a_1,\dots,a_n)$ 表示以 $a_1,\dots,a_n$ 为对角元、其余元素为零的 $n\times n$ 对角矩阵。

% ============================================================
\subsection{5.2 凸集}

\textbf{定义 5.1（凸集）}
集合 $C \subseteq \mathbb{R}^n$ 称为凸集，若对任意 $x,y \in C$ 和 $\lambda \in [0,1]$，有：
\[
\lambda x + (1-\lambda) y \in C
\]
几何意义：连接 $C$ 内任意两点的线段，整个都在 $C$ 内。"实心、无洞、无凹坑"。

\textbf{仿射集}（比凸集更强）：要求过任意两点的整条直线在集合内（而非仅线段）。仿射集 $\subseteq$ 凸集，线性方程组 $Ax=b$ 的解集是典型的仿射集。

（课件 p3.1, p3.2）

\textbf{常见凸集（正例）：}
\begin{itemize}
  \item \textbf{多面体} $C = \{x \mid Ax \le b\}$。线性不等式定义的区域，LP可行域 $\{x\mid Ax=b,\ x\ge 0\}$ 是它的特例。
  \item \textbf{单位圆盘} $C = \{(x_1,x_2) \mid x_1^2+x_2^2 \le 1\}$。由 $\ell_2$ 范数的三角不等式可证凸。
  \item \textbf{半空间} $C = \{x \mid a^T x \le b\}$。多面体的基本构件。
  \item \textbf{超平面} $\{x \mid a^T x = b\}$。既是仿射集也是凸集。
  \item \textbf{范数球} $\{x \mid \|x\| \le r\}$。任意范数定义的球都是凸集。
\end{itemize}

\textbf{凸集反例（为什么不是凸）：}
\begin{itemize}
  \item \textbf{圆环} $\{x \mid 1 \le x_1^2+x_2^2 \le 4\}$：取圆上对径两点，中点落在圆心空洞处——有洞。
  \item \textbf{坐标轴十字} $\{x \mid x_1 x_2 = 0,\ x_1 \ge 0,\ x_2 \ge 0\}$：$(1,0)$ 和 $(0,1)$ 的中点 $(0.5,0.5)$ 不在十字上——不实心。
  \item \textbf{整数点集} $\mathbb{Z}^n$：$0$ 和 $1$ 的中点 $0.5 \notin \mathbb{Z}$——离散不连通。
  \item \textbf{月牙形}：凹进去的部分使某些线段的中点落在集合外——有凹坑。
\end{itemize}

\textbf{直观判断法：} 脑子里想象集合形状——"实心无洞无凹坑"就是凸。

\subsubsection{凸组合与凸包}

\textbf{定义 5.2（凸组合与凸包）}
点 $x_1,\dots,x_k$ 的凸组合：$\sum_{i=1}^k \lambda_i x_i$，其中 $\lambda_i \ge 0$，$\sum \lambda_i = 1$。

集合 $S$ 的凸包 $\text{conv}(S)$：包含 $S$ 的最小凸集，等于 $S$ 中所有点的所有凸组合的集合。两点凸组合 $=$ 线段；三点凸组合 $=$ 三角形区域（含内部）。凸包就是把离散点集"填满"成包含它们的最小凸集。

（课件 p3.3）

\subsubsection{多面体}

\textbf{定义 5.3（多面体）}
多面体（polyhedral set）是有限个闭半空间的交：$\{x \mid Ax \le b\}$。LP的可行域 $\{x \mid Ax = b,\ x \ge 0\}$ 也是多面体——$Ax=b$ 等价于 $Ax\le b$ 且 $-Ax\le -b$，$x\ge0$ 等价于 $-x_i\le 0$。

（课件 p3.4, p3.5）

\subsubsection{极点}

\textbf{定义 5.4（极点）}
凸集 $C$ 中的点 $x$ 称为极点（extreme point），如果它不能表示为 $C$ 中两个不同点的严格凸组合。即：
\[
x = \lambda y + (1-\lambda)z,\ \lambda\in(0,1),\ y,z\in C \implies x=y=z
\]
极点是凸集的"角"——LP的基可行解（BFS）就是可行域多面体的极点。

（课件 p3.6）

\textbf{极点的几何直觉：}
\begin{itemize}
  \item 三角形区域的极点 = 三个顶点
  \item 正方形/长方体的极点 = 四个/八个顶点
  \item 圆盘 $\{x \mid \|x\|_2 \le 1\}$ 的极点 = 整个圆周（无穷多个极点）
  \item LP可行域 $\{x \mid Ax=b,\ x\ge0\}$ 的极点 = 基可行解（BFS）
\end{itemize}

\textbf{极点与单纯形法的关系：} LP的最优解一定能在极点（BFS）处取得。单纯形法就是在极点之间移动，每次迭代从一个BFS跳到目标函数值更优的相邻BFS。这是第2章单纯形法的理论基础。

\subsubsection{表示定理（选读）}

\textbf{定理 5.1（表示定理）}
设多面体 $P=\{x \mid Ax \le b\} \neq \emptyset$ 且 $\text{rank}(A)=n$。记 $P$ 的极点集为 $\{x^k\}_{k\in K}$（非空有限），极方向集为 $\{d^j\}_{j\in J}$（$J$ 可为空）。则：
\[
x \in P \;\Longleftrightarrow\; x = \sum_{k\in K}\lambda_k x^k + \sum_{j\in J}\mu_j d^j
\]
其中 $\lambda_k \ge 0,\ \sum\lambda_k = 1,\ \mu_j \ge 0$。此外，$J=\emptyset \Longleftrightarrow P$ 有界（称为多胞形）。

\textbf{推论：} 标准形多面体 $\{x \mid Ax=b,\ x\ge 0\}$ 非空则必有极点。这保证了单纯形法的理论基础——只需在极点（BFS）之间迭代搜索。

（课件 p3.8, p3.9）

% ============================================================
\subsection{5.3 凸函数}

\textbf{定义 5.5（凸函数）}
设 $C \subseteq \mathbb{R}^n$ 为非空凸集，$f: C \to \mathbb{R}$。$f$ 称为凸函数，若 $\forall x,y\in C,\ \lambda\in(0,1)$：
\[
f(\lambda x + (1-\lambda) y) \le \lambda f(x) + (1-\lambda) f(y)
\]
若 $x\neq y$ 时严格不等式成立，称为严格凸。若 $-f$ 是凸函数，则 $f$ 为凹函数；若 $-f$ 是严格凸函数，则 $f$ 为严格凹函数。

仿射函数（$a^Tx+b$）两边都取等号，既凸又凹。

（课件 p3.10）

\textbf{几何解释（弦在曲线上方）：} 任取定义域内两点 $x,y$，连接 $(x,f(x))$ 与 $(y,f(y))$ 的直线段整个在函数图像上方。凹函数则反过来——弦在图像下方。

\textbf{常见凸函数速查：}

\begin{tabular}{@{}lll@{}}
\toprule
\textbf{函数} & \textbf{定义域} & \textbf{凸性} \\
\midrule
$a^T x+b$（仿射） & $\mathbb{R}^n$ & 既凸又凹 \\
$x^2$    & $\mathbb{R}$   & 严格凸 \\
$e^{ax}$ & $\mathbb{R}$   & 凸 \\
$|x|$    & $\mathbb{R}$   & 凸（$x=0$不可微，但仍凸） \\
$x\ln x$ & $(0,\infty)$   & 凸（负熵） \\
$\|x\|$（任意范数）& $\mathbb{R}^n$ & 凸（三角不等式+齐次性）\\
$x_1^2+x_2^2$ & $\mathbb{R}^2$ & 严格凸 \\
$x_1^2-x_2^2$ & $\mathbb{R}^2$ & 非凸非凹（鞍面） \\
$x^3$    & $\mathbb{R}$   & 非凸非凹（$x>0$凸，$x<0$凹）\\
\bottomrule
\end{tabular}

% ============================================================
\subsection{5.4 凸函数的判别方法}

判断一个函数是否凸，有三层递进方法。三者的因果关系是逐步增加可微性假设，以此换取更强的计算可操作性：

\begin{center}
\fbox{%
\begin{minipage}{0.90\textwidth}
\small
\textbf{方法I：定义法}（无需可微假设，适合理论证明与简单函数）\\[2pt]
$\Big\Downarrow$ \quad 加上 $f$ 一阶可微\\[2pt]
\textbf{方法II：一阶条件/梯度不等式}（适合理论推导，但仍需"$\forall x,y$"）\\[2pt]
$\Big\Downarrow$ \quad 加上 $f$ 二阶可微\\[2pt]
\textbf{方法III：二阶条件/Hessian半正定}（最适合具体计算——顺序主子式法系统可操作）
\end{minipage}%
}
\end{center}

\textbf{为什么二阶判别最实用？} 定义法须证不等式对任意两点成立；一阶条件虽简化了形式，仍需在"所有 $x,y\in S$"上验证。二阶条件将凸性判定转化为矩阵半正定性的代数判定——计算Hessian矩阵，检查顺序主子式。对给定的具体函数，只需微积分 $+$ 初等代数，无需处理全称量词。这是考试中最常用的判别手段。

\subsubsection{方法I：定义法}

直接验证 $f(\lambda x+(1-\lambda)y) \le \lambda f(x)+(1-\lambda)f(y)$。无需可微性，仅适用于简单函数或理论证明。范数的凸性（不可微）只能用定义法证明——关键是用三角不等式和齐次性。

\subsubsection{方法II：一阶条件/梯度不等式}

设 $f$ 在开凸集 $S$ 上可微。则：
\[
f \text{ 凸 } \;\Longleftrightarrow\;
f(y) \ge f(x) + \nabla f(x)^T(y-x),\quad \forall x,y\in S
\]
几何含义：函数图像在任意点处的切线（一阶Taylor展开）始终在函数图像下方。

严格凸对应的不等号为严格 $>$（$\forall x\neq y$）。

（课件 p3.12）

\subsubsection{方法III：二阶条件/Hessian半正定（考试最常用）}

设 $f$ 在开凸集 $S$ 上二次可微。则：
\[
f \text{ 凸 } \;\Longleftrightarrow\; \nabla^2 f(x) \succeq 0,\quad \forall x\in S
\]
$\nabla^2 f(x) \succ 0$ 对所有 $x$ 成立 $\Longrightarrow$ $f$ 严格凸（充分但不必要！）

对于二次函数 $f(x)=\frac12 x^T Q x + c^T x$：$\nabla^2 f(x) \equiv Q$（常数矩阵），$f$ 凸 $\Longleftrightarrow$ $Q \succeq 0$。

（课件 p3.13）

\subsubsection{顺序主子式法——Hessian正/半正定的代数判据}

顺序主子式（leading principal minors）的定义：对 $n\times n$ 矩阵 $H$，第 $k$ 阶顺序主子式 $\Delta_k$ 是 $H$ 左上角 $k\times k$ 子矩阵的行列式：
\[
\begin{aligned}
\Delta_1 &= h_{11} \\
\Delta_2 &= \det\begin{pmatrix}h_{11}&h_{12}\\h_{21}&h_{22}\end{pmatrix} \\
\Delta_3 &= \det\begin{pmatrix}h_{11}&h_{12}&h_{13}\\h_{21}&h_{22}&h_{23}\\h_{31}&h_{32}&h_{33}\end{pmatrix}
\end{aligned}
\qquad\dots
\]

对 $n\times n$ 实对称矩阵 $H$（考试常见 $n=2$）：
\begin{itemize}
  \item $H \succ 0$（正定） $\Longleftrightarrow$ 所有顺序主子式 $> 0$（Sylvester判别法，充要）
  \item $H \succeq 0$（半正定）的判断需额外小心：顺序主子式 $\ge 0$ 是必要条件但不充分！

  \textbf{陷阱反例：} $A=\begin{pmatrix}0&0\\0&-1\end{pmatrix}$。$\Delta_1=0\ge0$，$\Delta_2=0\ge0$，两个顺序主子式都 $\ge0$。但取 $x=(0,1)^T$：$x^T A x = -1 < 0$，$A$ 不是半正定。原因是Sylvester判别法要求检查所有主子式（不只是顺序主子式）才能判定半正定。

  对常见 $2\times2$ 情形，实用充要条件是：$a\ge0,\ c\ge0,\ ac-b^2\ge0$ 三者同时成立 $\Longleftrightarrow$ $H\succeq0$。

  \item $2\times2$ 快速判据：$H=\begin{pmatrix}a&b\\b&c\end{pmatrix}$。
    \begin{itemize}
      \item $H\succ0 \Longleftrightarrow a>0,\ ac-b^2>0$
      \item $H\succeq0 \Longleftrightarrow a\ge0,\ c\ge0,\ ac-b^2\ge0$
      \item $H\prec0 \Longleftrightarrow a<0,\ ac-b^2>0$
      \item $H$ 不定 $\Longleftrightarrow ac-b^2<0$
    \end{itemize}
\end{itemize}

\subsubsection{严格凸 $\neq$ Hessian处处正定（重要陷阱）}

$\nabla^2 f \succ 0$（处处正定）是严格凸的充分但不必要条件。反例：$f(x)=x^4$（$\mathbb{R}$上），$f''(x)=12x^2 \ge 0$，$f''(0)=0$（非正定），但 $f$ 仍为严格凸函数。考试中若 Hessian 处处正定则直接下结论"严格凸"；若 Hessian 半正定但非正定，仍需结合其他方法确认是否严格凸。

% ============================================================
\subsection{5.5 凸规划}

\textbf{定义 5.6（凸规划）}
最优化问题 $\min\{f(x) \mid x \in S\}$ 称为凸规划，如果：
\begin{itemize}
  \item 可行域 $S$ 是凸集
  \item 目标函数 $f$ 是凸函数（对 $\min$ 问题；对 $\max$ 则要求 $f$ 凹）
\end{itemize}
两个条件缺一不可。

对于约束形式 $\min\{f(x)\mid g_i(x)\le 0,\ h_j(x)=0\}$，凸规划的条件等价于：
$f$ 凸，每个 $g_i$ 凸（不等式约束 $\le 0$ 形式），每个 $h_j$ 仿射（等式约束 $=0$ 形式）。

\textbf{定理 5.2（凸规划的核心性质）}
对于凸规划 $\min\{f(x) \mid x \in S\}$：
\begin{enumerate}
  \item 任意局部极小点都是全局极小点
  \item 若 $f$ 严格凸，则全局极小点唯一
  \item 上述结论不需要任何约束规格——仅需 $f$ 凸且 $S$ 凸
  \item 若约束规格成立，KKT条件（见第6章）升级为全局最优性的充要条件
  \item 全局极小点的集合是凸集
\end{enumerate}

（课件 p3.14）

\textbf{这为什么重要？} 第6章的KKT条件只是一阶必要条件——满足KKT的点不一定是全局最优。但一旦确认问题是凸规划，KKT点自动升级为全局最优解。这是连接第5章（凸性）和第6章（KKT）的关键桥梁。

\textbf{判定流程：} 目标函数凸？$\checkmark$ 可行域凸？$\checkmark$ $\Rightarrow$ 凸规划 $\Rightarrow$ KKT点=全局最优解。

\textbf{判定凸规划的完整操作步骤：}
\begin{itemize}
  \item Step 1：检查目标函数——求Hessian、算顺序主子式，判断是否半正定/正定。
  \item Step 2：检查可行域——是否可写为 $\{x\mid g_i(x)\le0,\ h_j(x)=0\}$ 形式，每个 $g_i$ 是否凸（不等号 $\le$）、每个 $h_j$ 是否仿射（等号）。
  \item Step 3：两个条件都满足 $\Rightarrow$ 凸规划 $\Rightarrow$ 局部$=$全局，KKT充要。任一不满足 $\Rightarrow$ 非凸规划 $\Rightarrow$ KKT仅为必要条件。
\end{itemize}

% ============================================================
\subsection{5.6 凸函数的保凸运算}

记住以下运算规则，可在不计算Hessian的情况下快速判断组合函数的凸性：

\begin{tabular}{@{}ll@{}}
\toprule
\textbf{运算} & \textbf{条件与结论} \\
\midrule
非负加权和  & $f,g$ 凸，$\alpha,\beta \ge 0$ $\Rightarrow$ $\alpha f+\beta g$ 凸 \\
与仿射复合  & $f$ 凸 $\Rightarrow$ $f(Ax+b)$ 凸 \\
逐点最大值  & $f_1,\dots,f_m$ 凸 $\Rightarrow$ $\max_i f_i(x)$ 凸 \\
逐点上确界  & $\forall y\in A$，$f(\cdot,y)$ 凸 $\Rightarrow$ $\sup_{y\in A} f(x,y)$ 凸 \\
标量复合(1) & $g$ 凸，$h$ 凸且单调不减 $\Rightarrow$ $h(g(x))$ 凸 \\
标量复合(2) & $g$ 凹，$h$ 凸且单调不增 $\Rightarrow$ $h(g(x))$ 凸 \\
\bottomrule
\end{tabular}

\textbf{实战技巧：} 看到 $\max$ 结构立刻想到保凸运算。例如 $f(x)=\max\{x_1^2+x_2^2,\; 2x_1+3x_2\}$——$x_1^2+x_2^2$ 凸（$\nabla^2 f \succ 0$），$2x_1+3x_2$ 凸（仿射），逐点最大值保凸 $\Rightarrow$ $f$ 凸。无需算Hessian。

\textbf{凸函数的下水平集：} 若 $f$ 为凸函数，则 $\{x \mid f(x) \le \alpha\}$ 是凸集（对任意 $\alpha$）。这为判断可行域凸性提供了快捷方法——例如 $\ell_1$ 范数球 $\{x\mid \|x\|_1\le 1\}$ 是凸集，因为范数是凸函数。

% ============================================================
\subsection{本章速查}

\begin{itemize}
  \item \textbf{凸集}：线段全在集合内。多面体（$Ax\le b$）、半空间、超平面、范数球都凸
  \item \textbf{凸集反例检验}：有洞（圆环）、有凹（月牙）、离散（$\mathbb{Z}^n$）、不连通状
  \item \textbf{仿射集} $\subseteq$ 凸集：要求过任意两点的整条直线在集合内
  \item \textbf{凸组合}：$\sum\lambda_i x_i$，$\lambda_i\ge0$，$\sum\lambda_i=1$
  \item \textbf{凸包} $\text{conv}(S)$：包含 $S$ 的最小凸集
  \item \textbf{多面体}：有限个闭半空间的交 $\{x\mid Ax\le b\}$
  \item \textbf{极点}：不能写成两个不同点严格凸组合的点。LP的BFS就是极点
  \item \textbf{表示定理}：多面体中任一点 = 极点凸组合 $+$ 极方向非负组合；标准形LP可行域非空则BFS存在
  \item \textbf{凸函数判别三层次}：定义法 $\to$ 一阶梯度不等式 $\to$ 二阶Hessian半正定
  \item \textbf{二阶判别最实用}：算Hessian $\to$ 顺序主子式 $\to$ 正定/半正定/不定/负定
  \item \textbf{$2\times2$ Hessian快判}：$a>0,\ ac-b^2>0$ 为正定；$ac-b^2<0$ 为不定
  \item \textbf{凸规划}：凸可行域 $+$ 凸目标（对 $\min$）$\Rightarrow$ 局部 $=$ 全局，KKT充要
  \item \textbf{凸规划判定}：目标凸性和可行域凸性缺一不可，必须分别验证
  \item \textbf{严格凸} $\Rightarrow$ 最优解唯一（但 Hessian 正定仅仅是充分条件）
  \item \textbf{保凸运算}：非负加权和、与仿射复合、逐点最大值均保持凸性
  \item \textbf{凸函数下水平集}：$\{x\mid f(x)\le\alpha\}$ 为凸集（$f$ 凸）
  \item \textbf{常见凸函数}：仿射、$x^2$、$e^x$、$|x|$、范数、$x\ln x$（负熵）
  \item \textbf{常见非凸/非凹}：$x_1^2-x_2^2$（鞍面）、$x^3$（$\mathbb{R}$上）、$\sin x$
\end{itemize}
